\section{Related Work}
\label{sec:related}

\paragraph{Recurrence and copy mechanisms.}
CyGNet~\cite{zhu2021cygnet} introduced a copy mode that scores the objects that already answered the same $(s, r)$ in the past, next to a generation mode over the whole vocabulary, on the observation that many facts recur. CENET~\cite{xu2023cenet} learns whether a query follows its historical or its non-historical dependency and applies the decision as a mask at inference. TiRGN~\cite{li2022tirgn} pairs a local recurrent graph encoder with a global history encoder for repeated facts and a time-guided decoder with periodicity. Gastinger et al.~\cite{gastinger2024history} showed that a non-learned recurrency baseline is competitive with eleven published models on three of five benchmarks, and CountTRuCoLa~\cite{gastinger2026counttrucola} turned that observation into four rule types with confidence functions that combine recency and frequency. All of these define ``history'' as the past of the query pair $(s, r)$. The dyad-only stratum of \cref{sec:diag}, where the answer interacted with $s$ under another relation, lies outside that definition by construction.

\paragraph{Snapshot evolution.}
RE-NET~\cite{jin2020renet} and RE-GCN~\cite{li2021regcn} evolve entity representations through a sequence of snapshot graphs, with a relational GCN per snapshot and a recurrent unit across them; CEN~\cite{li2022cen} adds a length-aware CNN with curriculum learning and an online setting; DiMNet~\cite{dong2025dimnet} adds cross-time links at equal hop depth and a disentanglement of active and stable node features. HiSMatch~\cite{li2022hismatch} frames the task as matching the query's historical structure against each candidate's. Our structural intensity is this family and is claimed as nothing more than the backbone; trained alone it reaches RE-GCN's level on ICEWS18 (\cref{sec:analysis}).

\paragraph{Global history graphs.}
The strongest recent results come from models that build a query-specific graph from the whole history and run message passing over it. LogCL~\cite{chen2024logcl} samples the one-hop facts of the query subject and the one-hop facts of the objects that answered $(s, r)$ before, applies entity-aware attention, and regularises the fusion of local and global views with four contrast patterns. HisRES~\cite{zhang2024hisres} aggregates all historical facts of the query pair into a global relevance graph and merges it with a multi-granularity recent encoder; its own ablation attributes six MRR points on ICEWS18 to that graph. CID-TKG and CHE-TKG~\cite{lei2026cidtkg,lei2026chetkg} add a second graph of recent facts retrieved through temporal logical rules and align the two views contrastively. CognTKE~\cite{chen2025cogntke} runs a fast one-hop and a slow multi-hop reasoner over a relation digraph. \model{} reads the same facts as event streams rather than as a graph: the stream encoder sees the order and the types of a dyad's events, the path intensity composes a stream with a typed edge, and nothing is contextualised by a second GNN. Of the three strongest members of this family, none has released code; the comparison in \cref{sec:exp} says so explicitly.

\paragraph{Paths, rules and explanation.}
xERTE~\cite{han2021xerte} expands a query-centred subgraph with temporal attention; TITer~\cite{sun2021titer} walks it with a reinforcement-learning agent; DaeMon~\cite{dong2023daemon} carries a path memory between the query subject and every candidate without entity representations, and TiPNN~\cite{dong2024tipnn} reasons over query-aware temporal paths on a unified history graph. TLogic~\cite{liu2022tlogic} extracts temporal logical rules by temporal random walks. Our untrained typed-dyad counter behaves like TLogic's length-one rules and reaches its ICEWS14 score; the path intensity is a learned form of its length-two rules, with the rule body read from an event stream instead of a single past fact.

\paragraph{Point processes.}
Know-Evolve~\cite{trivedi2017knowevolve} modelled each fact as a multivariate point process whose intensity is a function of evolving entity embeddings. GHNN~\cite{han2020ghnn} extended the neural Hawkes process to graph sequences and argued that treating every possible link as its own event-type dimension is infeasible because the number of types explodes. GAttNHP~\cite{tian2026gattnhp} encodes subject--relation chains as continuous-time point processes with grouped Hawkes priors. \model{} keeps the point-process reading and resolves the combinatorial objection: a dyad is instantiated only when it has history, so the per-pair process costs a few dozen short sequences per query, and its kernel is a learned function of a sinusoidal time encoding rather than a decay.

\paragraph{Diffusion models.}
DiffuTKG~\cite{cai2024diffutkg} casts the prediction as denoising the target fact with a Transformer conditioned on the query's history; NADEx~\cite{nadex2026} adds batch-wise negative prototypes to the conditioning and FreqDiff~\cite{freqdiff2026} a spectral calibration. These are orthogonal to the evidence question we study and are included in the comparison with their own reported numbers.

\paragraph{Evaluation.}
Gastinger et al.~\cite{gastinger2023apples} documented that filter settings, prediction horizon and dataset versions move TKG forecasting results by several points and distort comparisons; the recurrency baseline~\cite{gastinger2024history}, the co-occurrence shortcut of Zhang et al.~\cite{zhang2026evolution}, strikingness-aware reweighting~\cite{strikingness2026} and the synthetic shift benchmark of \"Ozdemir et al.~\cite{shift2026} all say that aggregate metrics on these benchmarks are dominated by repetition. The stratified evaluation of \cref{sec:analysis} is the real-data counterpart of probing mechanisms separately, and it is why we re-ran the one strong competitor with code under our protocol instead of quoting it.

\paragraph{Different settings.}
Inductive reasoning over emerging entities~\cite{transfir2026,adatkg2026} and continual fine-tuning at test time~\cite{hitscl2026} report gains under their own protocols. Our cold strata are not the emerging-entity problem: both entities are known, they have simply never interacted.
